\section{Introduction: one tree, four consumers}\label{sec:intro}

QtRocket is an open-source model-rocket simulator: a Qt6 Widgets GUI and a Qt-free command-line
REPL built on a shared C++23 simulation core, with a standalone OpenGL viewer alongside them. In
that core, a rocket is exactly one thing --- a tree of physical components. A nose cone sits on a
body tube; fins and a motor mount attach to the tube; a motor seats in the mount. The tree is not a
rendering convenience or a UI affordance: it is the physical description from which every quantity
the simulator needs (mass, center of mass, inertia, aerodynamic coefficients, thrust) is derived,
and every view the user sees is drawn.

This document describes the architecture that carries that tree: a pure-leaf component class
(\code{Part}), the owned node that composes leaves into assemblies (\code{PartNode}), the single
mutation authority that owns the whole structure (\code{PartsModel}), and the bridge that presents
the finished rocket to the numerics (\code{RocketModel}). This section establishes the shape of the
design at survey depth --- what the tree is, who reads it, and which class answers for what --- and
closes with a map of where each topic is developed in full.

\subsection{A rocket is a tree}\label{sec:intro-tree}

The leaves of the tree are concrete components, defined under \srcf{core/model/parts/}:
\code{ConicalNoseCone}, \code{BodyTube}, \code{FinSet}, \code{HollowSphere}, and \code{Motor}, a
zero-length node that wraps the time-aware \code{MotorModel} so that burning propellant mass is
just another leaf in the composite sum (\cref{sec:motor}). Each leaf knows only itself: its own
mass in $\unit{kg}$, its per-unit-mass inertia tensor about its own CM, its closed-form silhouette
radii, and its Barrowman aerodynamic contribution (\cref{sec:leaf}). Units are SI throughout, and
every part shares one local frame convention --- \zfwd, the part occupying $z \in [-L, 0]$
(\cref{sec:philosophy}).

The edges of the tree carry \emph{placement intent}, not coordinates: a \code{StationLink}
(\srcf{core/model/parts/Placement.h}) names a fractional station on the parent, a fractional
station on the child, a seat kind, and a gap in metres --- ``this rim seats against that rim,''
``this tube nests in that bore.'' A single resolver turns intent into absolute geometry once per
structural edit, and an envelope sweep passes an overlap verdict on the result (\cref{sec:placement}). Because
stations are fractions of live length, resizing a part moves every seam that references it; nothing
downstream stores a stale coordinate.

Finally, the empty tree is a real state. \code{PartsModel::hasDesign()} is false until a design is
installed, and every consumer is written against that possibility; there is no placeholder part
standing in for ``nothing yet'' (\cref{sec:tree}).

\subsection{Four consumers, one representation}\label{sec:intro-consumers}

Four independent subsystems consume the tree, and none of them holds a private copy of it:

\begin{enumerate}
\item \textbf{The CLI editor} (\srcf{cli/Repl.cpp}) builds and inspects designs interactively:
  \code{newdesign} installs a root, \code{addpart} resolves a parent by id, name, or the literal
  \code{root} and calls the \code{attach} verb (\src{cli/Repl.cpp}{952}), \code{listparts} walks
  the tree in attachment order printing composite mass and CG, \code{checkdesign} reports the
  envelope sweep's overlap verdict, and \code{savedesign}/\code{loaddesign} round-trip the design through
  the versioned XML serializer (\cref{sec:serialization,sec:cli}).
\item \textbf{The GUI tree view} (\srcf{gui/RocketTreeView.h}) adapts the live tree to Qt:
  \code{RocketPartModel} is a \code{QAbstractItemModel} whose rows \emph{are} the tree's children
  spans, and which consumes the typed aboutTo/did event stream so that an attach or detach becomes
  a surgical row insert or remove rather than a model reset (\cref{sec:gui}). Edits made elsewhere
  in the GUI route through the same verbs --- the cannonball preset writes its mass via
  \code{setPartMass} (\src{gui/CannonballTab.cpp}{130}).
\item \textbf{The 3D visualizer} (\srcf{visualizer/RocketMesh.cpp}) meshes each part at its
  resolved pose. It reads the root node's cached \code{resolvedPlacements()}
  (\src{visualizer/RocketMesh.cpp}{433}) --- the same poses the composite mass pass integrates
  over --- and reads \code{placementDiagnostics()} to paint overlap offenders in an error color
  with a located message.
\item \textbf{The flight integrator} (\srcf{core/sim/Propagator.cpp}) never sees parts at all. It
  drives the \code{Propagatable} interface, evaluating
  \cppinline|getForces(...) / getMass(t)| at every solver stage
  (\src{core/sim/Propagator.cpp}{41}) and recording composite mass properties once per recorded
  step. \code{RocketModel} answers every one of those calls from the root node's composite
  queries, as \cref{lst:intro-massprops} shows (\cref{sec:flight}).
\end{enumerate}

\begin{listing}[htbp]
\begin{minted}{cpp}
void RocketModel::writeMassProperties(double t, StateData& st)
{
    // Via the gated accessors, so after burnout this reads the frozen cache. The three are mutually
    // consistent at this t.
    st.mass    = parts_.root()->compositeMass(t);
    st.cg      = parts_.root()->compositeCm(t);
    st.inertia = parts_.root()->compositeI(t);
}
\end{minted}
\caption{The model side of the sim bridge delegates every mass property to the root node's cached
composite queries (\srccap{core/model/RocketModel.cpp}{49}). The integrator flies whatever the
tree says --- there is no second bookkeeping of mass anywhere in the sim path.}
\label{lst:intro-massprops}
\end{listing}

\begin{keyinsight}[Why one representation matters]
The pose the visualizer draws a coupler at and the CM the integrator flies are derived from the
\emph{same} cached resolve on the \emph{same} node; the mass column in the GUI and the CLI's
\code{listparts} read the same leaf values through the same \code{const} accessors. Consumers
cannot disagree about what the rocket is, because no consumer owns a copy that could drift. Where
representations must genuinely differ in shape --- Qt model indexes, GL vertex buffers --- they are
throwaway projections rebuilt from the tree, never parallel sources of truth.
\end{keyinsight}

\begin{rejected}[Rejected --- per-consumer scene copies]
The conventional alternative gives each front end its own mirror: the viewer keeps a scene graph,
the GUI keeps a shadow item tree, and change-propagation code keeps them synchronized. Every such
mirror is a place for the rocket to silently become two rockets --- a missed notification is not an
error anyone sees, it is a viewer that draws a design the integrator is not flying. The design here
spends its effort on making the one tree cheap to read (cached composites, a cached resolve,
$O(1)$ id lookup) instead of on synchronizing copies, so the failure mode does not exist.
\end{rejected}

\subsection{Three roles and a bridge}\label{sec:intro-roles}

The tree is carried by three classes with deliberately unequal jobs, plus one bridge above them.
\Cref{fig:intro-architecture} places them among their consumers.

\textbf{\code{Part}} (\src{core/model/parts/Part.h}{34}) is a pure leaf: one component's mass,
geometry profile, and aero contribution, with no children, no parent, and no knowledge that trees
exist. Its mutating setters are private to the \code{PartsModel} seam, so a leaf cannot be edited
behind the caches' back. \Cref{sec:leaf} gives its full contract.

\textbf{\code{PartNode}} (\src{core/model/PartsModel.h}{42}) is one node of the owned tree: it
owns its \code{Part}, stores the incoming \code{StationLink}, owns its children, and caches the
derived quantities of its sub-tree --- composite mass, CM, and inertia behind a mass-delta gate,
and the resolved placements plus overlap verdict behind a placement gate. Any node is a valid
sub-assembly root, so the same queries that describe the whole rocket describe a stage or a pod.
\Cref{sec:tree} covers the structure, \cref{sec:caching} the two gates, and \cref{sec:placement}
the resolve itself.

\textbf{\code{PartsModel}} (\src{core/model/PartsModel.h}{133}) is the single owner and the only
mutation authority. It holds the root, an $O(1)$ id index so stale references fail safe, and the
motor slot with its at-most-one invariant (\cref{sec:motor}). Every structural edit is a typed
verb --- \code{attach}, \code{attachSubtree}, \code{detach}, \code{installRoot} --- and the
fallible ones report \cppinline|std::expected| errors rather than throwing on misuse. Every
post-attach edit below the structural level is a routed verb --- \code{setPartMass} for mass,
\code{setLink} to re-seat a seam, \code{mutatePart} for arbitrary leaf edits under a declared
hint --- that writes and then dirties exactly the cache chain the edit touched. Around every mutation it fires a typed aboutTo/did event
pair, which is precisely the shape Qt's row-operation contract requires (\cref{sec:gui}).

\textbf{\code{RocketModel}} (\src{core/model/RocketModel.h}{27}) sits above as the rocket-level
facade: it owns the \code{PartsModel}, implements \code{Propagatable} for the sim
(\cref{sec:flight}), carries the drag configuration, and wraps design installation in
\code{installDesign} so install-time policy (such as resetting a manual reference-area override)
lives in one place.

\begin{designrule}[One representation, many readers, one writer]
Reads are \code{const} everywhere --- \code{children()}, \code{find()}, the composite queries, the
resolved placements. The only write paths are \code{PartsModel} verbs. Consequently every consumer
observes every edit, no edit can skip cache invalidation, and adding a fifth consumer is a matter
of reading the same tree, not of joining a synchronization protocol.
\end{designrule}

\tikzfig{component-architecture}{One tree, four consumers. The front ends and the integrator
(top and right) reach the owned part tree only through \code{RocketModel}'s facade and
\code{PartsModel}'s verbs and \code{const} reads; the typed event stream (dashed, left) flows back
out to the GUI adapter. \code{PartNode} calls the pure placement functions of
\srcfcap{core/model/parts/Placement.h} only when a structural edit dirties its geometry
cache.}{fig:intro-architecture}

\subsection{Where each topic lives}\label{sec:intro-map}

The rest of the document develops each role in the order the dependencies stack.
\Cref{sec:philosophy} states the design rules the whole structure obeys --- one mutation
authority, linear ownership, gated derived state --- and fixes the geometric conventions.
\Cref{sec:leaf} specifies the \code{Part} leaf and its concrete types; \cref{sec:tree} the
\code{PartNode} tree and \code{PartsModel}'s verbs, identity model, and typed errors.
\Cref{sec:caching} details the two cache gates and why the ODE-hot path never re-resolves
geometry; \cref{sec:placement} the station-link placement model, the per-edge pose math, and the
envelope sweep. \Cref{sec:motor} covers the motor slot and its single-motor invariant;
\cref{sec:serialization} the versioned \code{.qrd} design format. Two walkthroughs then trace the
system end to end: a CLI editing session in \cref{sec:cli} and a full simulated flight in
\cref{sec:flight}. \Cref{sec:gui} describes the Qt adapter built on the event stream,
\cref{sec:tradeoffs} records the load-bearing trade-offs and rejected alternatives in one place,
and \cref{sec:verification} documents the test architecture that pins all of it down.
